% Generated draft: requires mathematical and source review.
\documentclass[11pt]{article}
\usepackage[margin=25mm]{geometry}
\usepackage{fontspec}
\setmainfont{Latin Modern Roman}
\usepackage{amsmath,amssymb,mathtools,unicode-math}
\setmathfont{Latin Modern Math}
\usepackage{enumitem,needspace,xurl,hyperref}
\setlength{\parindent}{0pt}
\setlength{\parskip}{5pt}
\setlength{\emergencystretch}{3em}
\newcommand{\N}{\mathbb{N}}
\newcommand{\Z}{\mathbb{Z}}
\newcommand{\Q}{\mathbb{Q}}
\newcommand{\R}{\mathbb{R}}
\newcommand{\C}{\mathbb{C}}
\newcommand{\F}{\mathbb{F}}
\newcommand{\A}{\mathbb{A}}
\newcommand{\PP}{\mathbb{P}}
\newcommand{\E}{\mathbb{E}}
\newcommand{\eps}{\varepsilon}
\newcommand{\dd}{\,\mathrm{d}}
\begin{document}
\section*{UnsolvedMath ID: 30006286}
\textbf{Source problem code:} catalogue entry
\subsection{Definitions and mathematical statement}
A conformal modulus records the conformal equivalence class of a surface. For a fixed topological type \(\tau\), let \(\mathcal M_\tau\) denote the space of conformal equivalence classes of surfaces of that type, with boundary and marking conventions to be specified. Two surfaces represent the same class when there is a conformal isomorphism between them respecting those conventions.

Suppose that a Brownian surface \(B_\tau\) of type \(\tau\) has a specified probability law \(\mathbb P_\tau\) and a specified associated conformal structure. Write \(\operatorname{mod}_\tau(B_\tau)\in\mathcal M_\tau\) for its conformal class. Once the requisite measurable structures and measurability are established, its random conformal moduli have distribution
\[
\mu_\tau=(\operatorname{mod}_\tau)_*\mathbb P_\tau,
\qquad
\mu_\tau(A)=\mathbb P_\tau\!\left(\operatorname{mod}_\tau(B_\tau)\in A\right)
\]
for each measurable \(A\subseteq\mathcal M_\tau\).

\textbf{Problem, conditional on specifying the model.} For every topological type \(\tau\) covered by the intended Brownian-surface model, derive the distribution \(\mu_\tau\) of its conformal moduli, extending the Brownian-annulus case described as rigorously derived in the supplied record.

\textbf{Unresolved qualifications.} The record does not specify the Brownian-surface probability law, its normalization or conditioning, the construction of its conformal structure, or the admissible topological types. In particular, compactness, orientability, boundary components, punctures, markings, and any prescribed area or boundary lengths are unspecified. Thus the displayed formulation identifies the requested distribution only conditionally; it does not select a particular model or assert the existence of the conformal structure or moduli law. The retained original passage states the annulus result but does not explicitly state the general-topology extension; that extension appears in the supplied curated statement.
\subsection{Short English statement}
Determine the probability distribution of the conformal equivalence class of a Brownian surface for each intended topological type, extending the annulus case. A precise model and its boundary, marking, and normalization conventions remain unspecified. The supplied record labels this direction open; this draft does not independently verify that status.
\subsection{Sources}
\begin{itemize}
\item Supplied local UnsolvedMath registry record, permanent ID \(30006286\), entitled \emph{Random Moduli of General Brownian Surfaces}; exact source problem code \texttt{catalogue entry}. Collection: Oberwolfach Reports Open Problems. The record gives proposed year 2025 and names no individual proposer.
\item \emph{Statistical Physics Out of Equilibrium: Quantitative Results and Universality} (2025), Oberwolfach Reports, pp.\ 1333--1388, DOI \texttt{10.4171/owr/2025/26}. Supplied source and literature-source link, labeled ``original OWR report'': \url{https://doi.org/10.4171/owr/2025/26}. The supplied original passage reports that the random modulus had been rigorously derived only for the Brownian annulus; no pinpoint page or speaker attribution is supplied for that passage.
\item Supplied literature reference, labeled ``Ang--Remy--Sun, moduli of annuli'': \url{https://arxiv.org/abs/2203.12398}. The record attributes the annulus background to this reference without supplying a theorem number or its precise hypotheses.
\item Literature assessment dated exactly \texttt{2026-08-22}: the supplied curation describes the Brownian-annulus modulus as rigorously derived and the extension to general topology as a conjectural direction rather than a completed theorem. This is supplied dated curation, not an independently verified literature assessment. The background expressly requests independent verification of time-sensitive status and cited claims.
\item The registry labels the problem \texttt{open} and its statement \texttt{corrected\_verified}. Its supplied verification note says that the official report was checked and that the question was isolated and its notation repaired. These are recorded curation claims; no independent source visit or mathematical verification is asserted here.
\end{itemize}
\end{document}
